% Embedded Xinu Raspberry Pi 5 ユーザーズマニュアル
% コンパイル: lualatex xinu-pi5-manual.tex  (LuaTeX-ja / ltjsarticle)
\documentclass[a4paper,11pt]{ltjsarticle}

\usepackage{luatexja}
\usepackage[margin=24mm]{geometry}
\usepackage{listings}
\usepackage{xcolor}
\usepackage{enumitem}
\usepackage{longtable}
\usepackage[hidelinks]{hyperref}

\definecolor{codebg}{gray}{0.96}
\lstset{
  basicstyle=\ttfamily\small,
  backgroundcolor=\color{codebg},
  frame=single,
  rulecolor=\color{gray!40},
  breaklines=true,
  columns=fullflexible,
  keepspaces=true,
  showstringspaces=false,
  xleftmargin=4pt, xrightmargin=4pt,
  aboveskip=6pt, belowskip=6pt,
}
% コマンド片用ショートカット
\newcommand{\cmd}[1]{\texttt{#1}}

\title{\textbf{Embedded Xinu\\ Raspberry Pi 5 (BCM2712) \\ ユーザーズマニュアル}}
\author{Xinu-rpi5 プロジェクト}
\date{2026年6月}

\begin{document}
\maketitle
\tableofcontents
\newpage

%======================================================================
\section{はじめに}
本書は、Raspberry Pi 5（BCM2712 / Cortex-A76）上でベアメタル動作する
Embedded Xinu カーネル（\texttt{xinu-rpi5}）の使い方をまとめたユーザーズ
マニュアルである。本システムは MMU を有効化した AArch64 のシングルイメージ
カーネルで、以下を備える。

\begin{itemize}[noitemsep]
  \item HDMI フレームバッファとウィンドウシステム（デスクトップ）
  \item RP1 経由の有線イーサネット（TCP/IP・HTTP ゲートウェイ）
  \item 内蔵 CYW43455 による WiFi（スキャン／WPA2／DHCP／ping）
  \item RP1 xHCI による USB マウス・キーボード
  \item microSD カードの読み書き（FAT32）と \texttt{/microsd} マウント
  \item RAM 上で別カーネルへ切り替える \texttt{kexec} セレクタ
  \item 1 つのソースから派生する 4 種類の OS バリアント
\end{itemize}

\paragraph{重要な前提（起動メディア）.}
本機は\textbf{USB メモリから起動}し、本体の microSD スロットは通常空である。
USB メモリの FAT パーティション（\texttt{bootfs}）に置いた
\texttt{kernel\_2712.img} がファームウェアによって読み込まれて起動する。
microSD スロットは、本カーネルの SD ドライバが扱う\textbf{データ用}領域として
利用する（\texttt{/microsd}）。

%======================================================================
\section{システム構成}

\subsection{主要諸元}
\begin{longtable}{ll}
\hline
項目 & 値 \\
\hline
SoC & BCM2712 (Cortex-A76, AArch64) \\
カーネルイメージ & \texttt{kernel\_2712.img}（エントリ \texttt{0x80000}）\\
MMU & 有効（恒等マップ、キャッシュ OFF）\\
デバッグ UART & \texttt{0x107D001000}（3ピン JST、115200 8N1）\\
HDMI & 既定 1920$\times$1080$\times$32（OS3 は 1280$\times$720）\\
有線 IP（静的） & \texttt{192.168.3.101} \\
microSD コントローラ & sdhci-brcmstb sdio1（\texttt{0x10\_00FFF000}）\\
USB ホスト & RP1 内 DWC3/xHCI（PCIe 経由）\\
\hline
\end{longtable}

\subsection{コンソール}
出力は HDMI のテキストコンソールとデバッグ UART の両方に出る
（\texttt{uart\_putc} が画面にもミラーされる）。デスクトップ起動時は
ウィンドウシステムが画面を占有するため、テキストコンソールは隠れる。
ネットワーク経由のリモート操作（第\ref{sec:remote}節）が最も確実な操作手段である。

%======================================================================
\section{カーネルのフラッシュと起動}

\subsection{ビルド}
\texttt{compile/} ディレクトリで \texttt{make} を実行する。

\begin{lstlisting}[language=bash]
cd compile
make pi5        # フル OS        -> kernel_2712.img
make pi5-osmin  # OS1 (最小)     -> kernel_min.img
make pi5-os2    # OS2 (WMなし)   -> kernel_os2.img
make pi5-os3    # OS3 (低解像度) -> kernel_os3.img
\end{lstlisting}

\subsection{USB メモリへの書き込み（Mac）}
USB メモリ（FAT パーティション = \texttt{bootfs}）を Mac に挿し、
イメージをコピーして取り外す。マウントポイントを必ず確認すること。

\begin{lstlisting}[language=bash]
diskutil mount /dev/disk4s1
cp kernel_2712.img /Volumes/bootfs/kernel_2712.img
sync
diskutil unmount /dev/disk4s1
\end{lstlisting}

書き込んだら USB メモリを Pi 5 に戻し、電源を入れる。ファームウェアの
初期化に十数秒かかった後、起動ログが UART／HDMI に表示される。

%======================================================================
\section{4 つの OS バリアントと kexec による切り替え}
\label{sec:variants}

同一の \texttt{main.c} からコンパイルフラグだけを変えて、4 種類のカーネルを
生成できる。いずれも \texttt{kexec}（第\ref{sec:kexec}節）で相互に切り替え可能。

\begin{longtable}{lll}
\hline
名称 & イメージ & 特徴 \\
\hline
フル OS & \texttt{kernel\_2712.img} & ウィンドウシステム + ネットワーク（全機能）\\
OS1 & \texttt{OS1.IMG} & シェルのみ。ネットワークなし・ウィンドウなし \\
OS2 & \texttt{OS2.IMG} & 全機能だがウィンドウシステムなし（テキストコンソール）\\
OS3 & \texttt{OS3.IMG} & 全機能・HDMI を 1 ランク下げた 1280$\times$720 \\
\hline
\end{longtable}

\begin{itemize}[noitemsep]
  \item \textbf{OS1（最小）}：ネットワークもウィンドウも持たない。HDMI テキスト
        コンソールに \texttt{xinu-min\$} プロンプトを表示し、USB キーボードで操作する。
  \item \textbf{OS2（WM なし）}：ネットワーク等の全機能を持つが、ウィンドウ
        マネージャの代わりに全画面テキストコンソールでシェルを動かす
        （プロンプト \texttt{xinu-os2\$}）。ネットワーク経由の操作も可能。
  \item \textbf{OS3（低解像度）}：フル OS と同じだが HDMI を 1280$\times$720 に下げた版。
\end{itemize}

\paragraph{読みやすさ.} テキストコンソールのフォントは 8$\times$8 を 3 倍に
拡大して表示する（OS1・OS2 で有効）。

%======================================================================
\section{microSD ストレージ}
\label{sec:storage}

本体の microSD スロットは \texttt{/microsd} としてマウントされる
（カード未挿入でもディレクトリは存在する）。FAT32 カードの読み出しと、
小サイズファイルの永続書き込みに対応する。

\subsection{一覧と読み出し}
\begin{lstlisting}[language=bash]
ls /microsd                 # カード上のファイル一覧
cat /microsd/CONFIG.TXT     # ファイル内容の表示
\end{lstlisting}

\subsection{書き込み（永続）}
HTTP 経由で 1 クラスタ（32\,KB）以下の小ファイルをカードへ永続書き込みできる。
ボディが内容、URL 末尾がファイル名（8.3 形式）。

\begin{lstlisting}[language=bash]
curl -X POST --data-binary "hello from xinu" \
     http://192.168.3.101/microsd/write/TEST.TXT
# 再起動後も ls /microsd / cat /microsd/TEST.TXT で残存
\end{lstlisting}

\subsection{診断}
SD コントローラとカードの初期化状況を表示する診断コマンド。
\begin{lstlisting}[language=bash]
sdtest        # コントローラ／CMD0,CMD8,ACMD41 等の応答を表示
\end{lstlisting}

%======================================================================
\section{kexec カーネルセレクタ}
\label{sec:kexec}

\texttt{kexec} は microSD 上のカーネルイメージを RAM に読み込み、
電源を切らずにそのカーネルへジャンプする（再フラッシュ不要）。失敗しても
RAM 上の操作なので、電源再投入で USB のフル OS に戻るだけでブリック риск はない。

\begin{lstlisting}[language=bash]
kexec /microsd/OS1.IMG      # OS1 (最小) を起動
kexec /microsd/OS2.IMG      # OS2 (WMなし) を起動
kexec /microsd/OS3.IMG      # OS3 (低解像度) を起動
kexec /microsd/KERNEL~1.IMG # カード上のフル OS へ
\end{lstlisting}

\paragraph{注意点.}
\begin{itemize}[noitemsep]
  \item \texttt{kexec} の後はネットワークが落ちる（非ファームウェア状態からの
        PCIe 再初期化のため）。結果は HDMI 画面で確認する。
  \item ブートできるのは\textbf{ベアメタルの Xinu カーネルのみ}。Linux カーネル
        （例 \texttt{KERNEL8.IMG}、約 9.6\,MB）は別のブートプロトコルを必要とし、
        単純なチェインロードではハングする。本実装は 4\,MB 以上のイメージを
        Linux と見なして拒否する（Xinu と Linux はどちらも arm64 Image で
        マジックが同じため、サイズで判別する）。
\end{itemize}

%======================================================================
\section{シェルコマンド}
画面のシェル（USB キーボード）か、ネットワーク経由（第\ref{sec:remote}節）の
\texttt{/run?cmd=...} から実行する。主なコマンドを示す。

\begin{longtable}{ll}
\hline
コマンド & 説明 \\
\hline
\cmd{help} & コマンド一覧 \\
\cmd{ls [path]} & ディレクトリ一覧 \\
\cmd{cat <file>} & ファイル内容表示 \\
\cmd{cd / pwd} & カレントディレクトリ移動／表示 \\
\cmd{kexec /microsd/<k.img>} & 別カーネルを起動（第\ref{sec:kexec}節）\\
\cmd{cc <file> / make <file>} & C/AIPL プログラムをその場でコンパイル・実行 \\
\cmd{mem} & ヒープ／BSS の状況 \\
\cmd{peek <hex>} & 32bit MMIO ワードの読み出し \\
\cmd{uptime} & 汎用タイマ値 \\
\cmd{ps} & コア／EL 状態 \\
\cmd{wifi on/off/status/scan} & 接続／切断／状態／スキャン（第\ref{sec:wifi}節）\\
\cmd{wifi adhoc <ssid> [ch] [n]} & MANET ad-hoc（メッシュ）セルに参加（第\ref{sec:wifi}節）\\
\cmd{wifi aodv <a.b.c.d>} & マルチホップのメッシュ経路を探索（第\ref{sec:wifi}節）\\
\cmd{sdtest} & SD コントローラ診断 \\
\cmd{reboot} & 再起動（PM ウォッチドッグ）\\
\cmd{clear} & 画面クリア \\
\hline
\end{longtable}

%======================================================================
\section{ネットワーク経由のリモート操作}
\label{sec:remote}
フル OS・OS2・OS3 は有線 IP \texttt{192.168.3.101} で HTTP を待ち受ける。
シリアル入力が使いにくい環境でも、Mac 等から \texttt{curl} で操作できる。

\begin{longtable}{ll}
\hline
エンドポイント & 説明 \\
\hline
\cmd{/run?cmd=<shell>} & 任意のシェルコマンドを実行し出力を返す \\
\cmd{/fs}, \cmd{/fs/ls/<d>}, \cmd{/fs/cat/<f>} & VFS ツリーの参照 \\
\cmd{/fs/write/<f>} (POST) & VFS（tmpfs）へファイル書き込み \\
\cmd{/microsd/write/<NAME>} (POST) & microSD へ永続書き込み（第\ref{sec:storage}節）\\
\cmd{/usb/...} & USB の列挙・診断（第\ref{sec:usb}節）\\
\cmd{/api/actors}, \cmd{/send?to=\&m=} & アクタ／AIPL ゲートウェイ \\
\hline
\end{longtable}

\begin{lstlisting}[language=bash]
curl 'http://192.168.3.101/run?cmd=ls%20/microsd'
curl 'http://192.168.3.101/run?cmd=cat%20/microsd/CONFIG.TXT'
curl 'http://192.168.3.101/run?cmd=wifi%20status'
\end{lstlisting}

\noindent\textbf{補足}：URL 中の空白は \texttt{\%20} もしくは \texttt{+} で
エンコードする。出力バッファには上限があるため、長い結果は分割される。

%======================================================================
\section{WiFi}
\label{sec:wifi}
内蔵 CYW43455 によるフル WiFi（スキャン／WPA2／DHCP／ping）に対応する。
\textbf{起動時には自動接続しない}。\cmd{wifi on} を実行したときのみ接続する。

\subsection{アクセスポイントへの接続}
\begin{lstlisting}[language=bash]
wifi on <ssid> <pass>   # ファームウェア起動(~10s) -> WPA2 join -> DHCP
wifi on                 # 直前の資格情報、またはビルド時の wifi.conf を再利用
wifi scan               # 無線を起こして近傍 AP を一覧
wifi status             # 接続状態・SSID・IP
wifi off                # 無線停止
wifi-invest             # 接続できないときの診断
\end{lstlisting}
\cmd{wifi on} 成功時は \texttt{wifi: CONNECTED IP=...} で終わる。画面右下に WiFi の
電波アイコンと、その下に SSID・自 IP を表示する。資格情報はビルド時に
\texttt{.gitignore} 済みの \texttt{wifi.conf} からカーネルへ埋め込まれる
（パスワードはコミットしない）。

\subsection{メッシュネットワーク（MANET ad-hoc）}
複数の Pi 5（および Pi 4 / Pi 3）Xinu ノードを\textbf{アクセスポイント無し}で直接
つなぎ、ピアツーピアのメッシュを組める。802.11 の IBSS（ad-hoc）モード +
オンデマンドの AODV マルチホップルーティングを使う。

\begin{lstlisting}[language=bash]
wifi adhoc <cell-ssid> [ch] [node]   # ad-hoc セルに 10.0.0.<node> で参加
wifi aodv  <a.b.c.d>                 # マルチホップ経路をオンデマンド探索
\end{lstlisting}

\begin{itemize}[noitemsep]
  \item \cmd{wifi adhoc mesh1 6 1} は指定セル（BSSID \texttt{02:4d:41:4e:45:54}
        = ``MANET''）のチャンネル 6（既定）で無線を IBSS モードで起こし、静的 IP
        \texttt{10.0.0.1}（\texttt{[node]} 番号、既定 1）を取る。以降 AODV 中継が
        自動で動く。
  \item 全ノードで\textbf{同じセル SSID・同じチャンネルを使い、\texttt{[node]}
        番号は重複させない}（\texttt{10.0.0.<node>/24}）。電波の届く範囲のノードは
        \texttt{10.0.0.x} で直接通信できる。
  \item 直接届かない宛先には \cmd{wifi aodv 10.0.0.3} で RREQ をブロードキャスト
        （UDP 654）し、中間ノードが中継、宛先が RREP を返してマルチホップ経路を
        確立する（例 \texttt{AODV route: 10.0.0.3 via 10.0.0.2, 2 hop}）。先に
        \cmd{wifi adhoc} で IP を取得しておくこと。
\end{itemize}
ad-hoc/AODV はインフラ（\cmd{wifi on}）モードとは独立で、再起動後は復元しないため、
各ノードで \cmd{wifi adhoc} を再実行する。

%======================================================================
\section{USB（マウス・キーボード）}
\label{sec:usb}
USB-A ポートは RP1 内の 2 つの DWC3/xHCI コントローラに接続される。
ブート時のフォールトを避けるため、列挙は起動数秒後に自動で行われる
（あるいは \texttt{/usb/...} エンドポイントから手動駆動）。
ブートマウスはカーソル移動に、キーボードはシェル入力に用いられる。

%======================================================================
\section{付録}

\subsection{ビルドターゲット一覧}
\begin{longtable}{ll}
\hline
ターゲット & 生成物 \\
\hline
\cmd{make pi5} & \texttt{kernel\_2712.img}（フル OS）\\
\cmd{make pi5-osmin} & \texttt{kernel\_min.img}（OS1）\\
\cmd{make pi5-os2} & \texttt{kernel\_os2.img}（OS2）\\
\cmd{make pi5-os3} & \texttt{kernel\_os3.img}（OS3）\\
\hline
\end{longtable}

\subsection{トラブルシューティング}
\begin{itemize}[noitemsep]
  \item \textbf{HTTP が応答しない}：起動直後はネットワーク初期化前。十数秒
        待ってから再試行する。\texttt{kexec} 後はネットワークが落ちるため、
        電源再投入で USB のフル OS に戻る。
  \item \textbf{\texttt{kexec} で「too large」と出る}：4\,MB 以上のイメージ
        （Linux カーネル等）はチェインロードできない。ベアメタル Xinu の
        イメージを指定すること。
  \item \textbf{\texttt{/microsd} が空}：スロットに microSD が入っているか確認。
        \texttt{sdtest} で初期化状況を確認できる。
  \item \textbf{画面が出ない}：HDMI フレームバッファのメールボックス初期化に
        失敗した場合、コンソールは UART のみとなる。
\end{itemize}

\subsection{安全上の注意}
\begin{itemize}[noitemsep]
  \item WiFi のパスワードはソースリポジトリにコミットしない（ビルド時に
        \texttt{.gitignore} 済みファイルから埋め込む運用）。
  \item \texttt{kexec} は RAM 上の操作であり、不正なイメージでも電源再投入で
        復帰する（フラッシュ破壊のリスクはない）。
\end{itemize}

\vfill
\begin{center}\small Embedded Xinu / Raspberry Pi 5 — 2026年6月\end{center}

\end{document}
